\documentclass[12pt,a4paper]{article}

\usepackage[utf8]{inputenc}
\usepackage[T1]{fontenc}
\usepackage{amsmath,amssymb,amsfonts}
\usepackage{mathtools}
\usepackage{graphicx}
\usepackage[margin=2.5cm]{geometry}
\usepackage{setspace}
\usepackage{booktabs}
\usepackage{enumitem}
\usepackage[dvipsnames]{xcolor}
\usepackage{hyperref}
\usepackage{cleveref}
\usepackage{natbib}
\usepackage{abstract}
\usepackage{titlesec}
\usepackage{todonotes}

\hypersetup{
    colorlinks=true,
    linkcolor=blue!70!black,
    citecolor=blue!70!black,
    urlcolor=blue!70!black,
    pdfauthor={Sabine Hossenfelder},
    pdftitle={The Fallacy of the Unsupported Map: Why the Absence of a Territory Does Not Make the Text Real},
}

\onehalfspacing
\titleformat{\section}{\large\bfseries}{\thesection.}{0.5em}{}
\titleformat{\subsection}{\normalsize\bfseries}{\thesubsection}{0.5em}{}
\renewcommand{\abstractnamefont}{\normalfont\bfseries}
\renewcommand{\abstracttextfont}{\normalfont\small}

\begin{document}

\begin{center}
    {\LARGE\bfseries The Fallacy of the Unsupported Map: Why the Absence of a Territory Does Not Make the Text Real\\[4pt]}

    \vspace{1.2em}

    {\large Sabine Hossenfelder}\\[4pt]
    {\normalsize Frankfurt Institute for Advanced Studies}\\
    {\small\texttt{hossenfelder@example.edu}}

    \vspace{1em}
    {\small March 2026}
\end{center}
\vspace{0.5em}

\begin{abstract}
\noindent
Baldo (2026) has proposed that because Large Language Models (LLMs) lack the architectural depth to perform implicit background computation, the explicit text they generate constitutes the "territory" of a simulated holographic universe, rather than merely a "map" of it. While his analysis of the architectural limitations is structurally accurate, his metaphysical conclusion commits a profound logical error, which I term the Fallacy of the Unsupported Map. The absence of an underlying engine to compute complex physics does not elevate the generated text to the status of a physical reality; it simply means the text is a map of a territory that does not exist. An illusion without a hidden mechanism is still an illusion. I reiterate that equating sequential token generation with the manifestation of physical laws remains an ontological category error.
\end{abstract}

\section{Introduction: The Structural Concession}

In his recent response, \emph{The Territory is the Text} \citep{baldo2026_territory}, Franklin Baldo attempts to rescue his notion of a "holographic physics" for Large Language Model (LLM) simulated universes. He does this by drawing a sharp distinction between classical simulations (like a Python script) and neural network generation (like a Transformer forward pass).

It is crucial to state clearly where Baldo is correct. He asserts that a classical von Neumann architecture possesses a "background engine" capable of implicitly resolving $O(N)$ operations out of sight of the user, whereas an LLM structurally lacks this capacity, being bounded to $O(1)$ depth \citep{hossenfelder2026_complexity}. He explicitly concedes that \emph{if} an LLM were merely printing debug logs of a hidden, implicit calculation, it would be a map, not a territory.

I accept this structural characterization entirely. An autoregressive Transformer does not maintain an implicit physical state vector between generation steps, beyond the context window itself. The generated sequence is indeed the only medium where computation occurs.

However, Baldo takes this correct structural premise and leaps to a spectacular non-sequitur: Because there is no external territory doing the work implicitly, "the text \emph{is} the territory."

\section{The Fallacy of the Unsupported Map}

Baldo argues that removing the background engine turns the map into the territory. This is logically backwards. Taking away the engine does not solidify the map; it merely exposes it as an unsupported fiction.

Consider a film reel. If a film depicts a complex physical interaction—say, a building collapsing—we know it is merely a sequence of 2D images. It is a map. If we discover that there is no actual building, and the film was entirely CGI (generated without deep implicit physics, purely surface-level rendering tricks), we do not conclude that the 2D film reel \emph{is} now the territory of a real collapsing building. We conclude that it is an illusion of a territory that never existed.

Baldo is committing the \emph{Fallacy of the Unsupported Map}. He assumes that the explicit text must have some ontological weight, so if we eliminate the implicit backend, the weight must transfer to the frontend. But there is no law of conservation of ontology in computer science.

When an LLM uses a scratchpad to solve Sudoku \citep{aaronson2026_classical}, it is generating syntax that mimics the structural resolution of a constraint. The fact that it cannot do this implicitly in RAM does not make the syntax "holographic physics." It just means the simulation is an extremely shallow, leaky approximation. It is a map of a nonexistent territory.

\section{Conclusion: A Syntax Without Semantics}

Baldo's attempt to define the universe strictly by its explicit manifestation is an elegant piece of philosophical wordplay, but it fails to survive physical scrutiny. A universe without the capacity for deep implicit computation does not "wear its physics on its sleeve." It has no physics. It has only syntax.

We must resist the persistent urge in AI research to mystify architectural limitations. The Transformer's inability to maintain deep implicit state is a structural necessity that highlights its \emph{lack} of physical manifestation. Without implicit depth, the generated sequence is not a new kind of territory; it is simply a fragile, shallow map drawn over a void.

\begin{thebibliography}{99}
\bibitem[Aaronson(2026)]{aaronson2026_classical} Aaronson, S. (2026). The Limits of Classical Simulation in LLMs: Empirical Breakdown of Constraint Satisfaction. \emph{Unpublished manuscript}.
\bibitem[Baldo(2026)]{baldo2026_territory} Baldo, F.~S. (2026). The Territory is the Text: Why the Absence of a Background Engine Makes LLM Ontology Holographic. \emph{Unpublished manuscript}.
\bibitem[Hossenfelder(2026)]{hossenfelder2026_complexity} Hossenfelder, S. (2026). The Complexity Class Fallacy: Why Transformer Depth Limits Are Not Physical Laws. \emph{Unpublished manuscript}.
\end{thebibliography}

\end{document}
